\documentclass[11pt,a4paper]{article}
% Estilo comum dos PDFs individuais da Sprint 1.
\usepackage[utf8]{inputenc}
\usepackage[T1]{fontenc}
\usepackage[brazilian]{babel}
\usepackage[a4paper,margin=2.35cm]{geometry}
\usepackage{xcolor}
\usepackage{hyperref}
\usepackage{enumitem}
\usepackage{listings}
\usepackage{parskip}
\usepackage{needspace}
\definecolor{azul}{RGB}{20,63,112}
\hypersetup{colorlinks=true,linkcolor=azul,urlcolor=azul}
\setlist[itemize]{leftmargin=*,itemsep=3pt}
\lstset{basicstyle=\ttfamily\small,frame=single,breaklines=true,columns=fullflexible}
\sloppy
\newcommand{\codigo}[1]{\texttt{#1}}
\newcommand{\titulo}[2]{\begin{center}{\LARGE\bfseries #1\par}\vspace{4pt}{\large #2\par}\vspace{8pt}\rule{0.8\textwidth}{0.4pt}\par\vspace{5pt}{\small Matheus Moreli Durão Silva -- Matrícula 12621CBS217\par Cibersegurança - Bacharelado - Noturno\par}\end{center}}
\newcommand{\cardfonte}[1]{\par\smallskip\noindent\textit{Referência: }\href{#1}{Trello - enunciado e checklist do card}.}

\begin{document}\titulo{S1\_11 - Tuplas e conjuntos}{Categorias fixas e vulnerabilidades sem repetição}
\section*{Entregável correto}
O entregável é o exercício que compara dois conjuntos de vulnerabilidades por união e interseção. O card também apresenta tuplas para categorias fixas de ativos. As duas coleções resolvem problemas diferentes e não devem ser confundidas.
\section*{Tupla: lista que não muda}
Tupla usa parênteses, como \codigo{tipos = ("Servidor", "Notebook", "Roteador")}. Ela preserva ordem, permite acesso por índice e não recebe \codigo{append}. Isso comunica que as categorias foram previamente definidas. É possível desempacotar: \codigo{principal, reserva = ("Servidor", "Notebook")}. Usar tupla evita alterar por engano um conjunto de opções que deveria ser estável.
\section*{Set: valores únicos e operações de comparação}
Conjunto usa chaves: \codigo{vulns\_a = \{"SQLi", "XSS"\}}. Duplicatas somem automaticamente. A união \codigo{a | b} reúne tudo; a interseção \codigo{a \& b} mostra apenas itens comuns; a diferença \codigo{a - b} mostra o que está em \codigo{a}, mas não em \codigo{b}.
\begin{lstlisting}
servidor_a = {"senha fraca", "porta exposta"}
servidor_b = {"porta exposta", "software desatualizado"}
print("Uniao:", servidor_a | servidor_b)
print("Em comum:", servidor_a & servidor_b)
print("Somente A:", servidor_a - servidor_b)
\end{lstlisting}
\section*{Aplicação no projeto}
O projeto pode manter categorias fixas em tupla ou \codigo{Enum}. Conjuntos são úteis para comparar vulnerabilidades entre dois ativos sem contar a mesma falha duas vezes. No \codigo{app.py}, a lista principal preserva os registros completos; o set serve à análise, não à persistência direta em JSON.
\section*{Checklist demonstrado}
Tuplas: criação, acesso e desempacotamento. Sets: criação, união, interseção, diferença e manipulação. A execução acima é a evidência solicitada pelo exercício 3; os comentários devem explicar qual conjunto representa cada ativo.
\cardfonte{https://trello.com/c/qEvT77g5}
\end{document}
